% !TEX root = ../main.tex
% ESEMPIO 0d: slide degli esperimenti e dei risultati.
%   - Experiments (1): testo a sinistra e immagine a destra in due colonne; \pause fa comparire il punto successivo.
%   - Experiments (2): intestazione fissa e contenuto che cambia con \only<n>; overlayarea riserva lo spazio,
%     così niente si sposta da un passo all'altro. Il terzo passo è una tabella (vedi anche ex_05_tables.tex).
%   - Results: grafico a barre orizzontali (pgfplots, groupplot) con dati letti da un CSV (vedi ex_04_line_plots.tex
%     per i grafici a linee).
%   Da copiare: il frame che serve, poi sostituire i CHANGE.

\begin{frame}
    \frametitle{Example: experiments, text and image}
    % GUIDA: obiettivo degli esperimenti e come sono impostati. Se sono tanti, mostra solo il sottoinsieme che fa passare l'intuizione.

    \alert{Goal}: CHANGE: goal of the experiments

    \vspace{1em}

    % Testo a sinistra, immagine a destra. \pause fa comparire il punto successivo al passo seguente.
    \begin{columns}
        \begin{column}{.7\textwidth}
            \begin{itemize}\setlength{\itemsep}{1em}
                \item \alert{CHANGE: item 1}: CHANGE: description of the first point\ldots
                      \pause\item \ldots CHANGE: item 2, which appears at the next step
            \end{itemize}
        \end{column}
        \begin{column}{.3\textwidth}
            \includegraphics[width=.6\textwidth]{assets/img/patrick_star_sad.png}
        \end{column}
    \end{columns}

\end{frame}

\begin{frame}
    \frametitle{Example: experiments, changing content}
    % GUIDA: come è costruito l'esperimento (dati, misure, configurazioni); la tabella riassume le configurazioni.

    % Intestazione fissa e contenuto che cambia: \only<1-> resta visibile a tutti i passi, mentre
    % \only<1>, \only<2>, \only<3> compaiono solo al passo indicato. overlayarea riserva lo spazio,
    % così niente si sposta da un passo all'altro.
    \begin{overlayarea}{\textwidth}{.75\textheight}
        \only<1->{
                \alert{CHANGE: fixed header}: CHANGE: sentence that stays at every step, with \alert{key words} and the symbol $T$

                \vspace{1em}
            }
        \only<1>{
                CHANGE: introduction of the list, in terms of

                \begin{columns}
                    \begin{column}{.8\textwidth}
                        \begin{itemize}\setlength{\itemsep}{1em}
                            \item CHANGE: point 1
                            \item CHANGE: point 2
                                  \begin{itemize}
                                      \item CHANGE: sub-point of point 2, with an \emph{emphasized} word
                                  \end{itemize}
                        \end{itemize}

                        \vspace{1em}

                        CHANGE: closing sentence
                    \end{column}
                    \begin{column}{.2\textwidth}
                        \includegraphics[width=1\textwidth]{assets/img/patrick_scientist.jpg}
                    \end{column}
                \end{columns}
            }
        \only<2>{

                \begin{columns}
                    \begin{column}{.8\textwidth}
                        CHANGE: sentence about the \alert{key concept} $T$

                        \vspace{1em}

                        \begin{itemize}\setlength{\itemsep}{1em}
                            \item CHANGE: point 1
                            \item CHANGE: point 2
                        \end{itemize}
                    \end{column}
                    \begin{column}{.2\textwidth}
                        \includegraphics[width=1\textwidth]{assets/img/patrick_question.jpg}
                    \end{column}
                \end{columns}
            }
        \only<3>{
                \vspace{1em}

                % Tabella con intestazione su due righe: \multirow per la prima colonna, \multicolumn per i gruppi.
                % Le righe sono di booktabs (\toprule, \midrule, \bottomrule); \cmidrule traccia una riga parziale.
                \begin{adjustbox}{max totalsize={.675\linewidth}{\textheight},center}
                    \scriptsize
                    \begin{tabular}{c|cc|cc|cc}
                        \toprule
                        \multirow{2}{*}{\textbf{CHANGE: row label}} & \multicolumn{2}{c|}{\textbf{CHANGE: group A}} & \multicolumn{2}{c|}{\textbf{CHANGE: group B}} & \multicolumn{2}{c}{\textbf{CHANGE: group C}}                                                                                                                                \\
                                                                    & \multicolumn{1}{c}{\textbf{CHANGE: A1}}       & \multicolumn{1}{c|}{\textbf{CHANGE: A2}}      & \multicolumn{1}{c}{\textbf{CHANGE: B1}}      & \multicolumn{1}{c|}{\textbf{CHANGE: B2}} & \multicolumn{1}{c}{\textbf{CHANGE: C1}} & \multicolumn{1}{c}{\textbf{CHANGE: C2}} \\
                        \midrule
                        1                                           & 0                                             & 0                                             & 0                                            & 0                                        & 4                                       & 16                                      \\
                        2                                           & 80                                            & 20                                            & 0                                            & 0                                        & 0                                       & 0                                       \\
                        3                                           & 10                                            & 20                                            & 15                                           & 5                                        & 5                                       & 5                                       \\
                        4                                           & 20                                            & 5                                             & 10                                           & 5                                        & 5                                       & 10                                      \\
                        \bottomrule
                    \end{tabular}
                \end{adjustbox}
            }
    \end{overlayarea}

\end{frame}

\begin{frame}
    \frametitle{Example: results plot}
    % GUIDA: i risultati, per esempio come grafici, con una sola conclusione per grafico.

    % Grafico a barre orizzontali (pgfplots, groupplot) con dati letti da un CSV; vedi il file incluso.
    % Due pannelli per slide. Per altre configurazioni o per la variante a quattro pannelli vedi i commenti in testa al file.
    % GRAFICO: due pannelli affiancati (una configurazione ciascuno), barre orizzontali, dati letti da un CSV,
% una sola legenda condivisa.
%
% Colonne di data.csv (x index): 0 e 1 = configurazione 1 (stato dell'arte, nuovo approccio); 2 e 3 = configurazione 2.
%
% VARIANTI
% - Altre configurazioni: il CSV ha anche le colonne 4–7 (configurazioni 3 e 4). Per mostrarle in una seconda slide
%   si copia questo file, si cambiano i titoli dei pannelli ("Configuration N") e gli x index (4, 5 e 6, 7).
% - Quattro pannelli (2x2): group size=2 by 2 e quattro \nextgroupplot. In una sola slide non entra (la riga di sotto esce
%   dalla pagina), quindi è adatto a una figura di un articolo. In quel caso, sotto ogni pannello si possono mettere
%   le lettere (a)-(d) con un ciclo dopo \end{groupplot}:
%       \foreach \letter/\target in {a/c1r1, b/c2r1}{ \node[below=0.25cm of my plots \target.south] {\small (\letter)}; }
%   dove cNrM è il pannello di colonna N e riga M del gruppo "my plots". Per la riga di sotto, che ha anche l'etichetta
%   dell'asse x, la distanza va portata a 0.8cm.
% - Legenda: è definita nel secondo pannello (legend to name=leg) e posizionata sopra i due pannelli con
%   \pgfplotslegendfromname. Lo stesso nome può essere definito una volta sola nel documento: se due slide includono
%   un grafico con questo codice, il secondo deve usare un nome diverso (per esempio leg2).
\pgfplotstableread[col sep=comma]{assets/tex/03_plot_quality/data.csv}\data

\begin{adjustbox}{max totalsize={.99\linewidth}{\textheight},center}
    \begin{tikzpicture}
        \pgfplotsset{every tick label/.append style={font=\normalsize}}
        \pgfplotsset{every axis label/.append style={font=\normalsize}}
        \begin{groupplot}[
                group style={
                    group name=my plots,
                    group size=2 by 1,
                    vertical sep=1.45cm, 
                    xticklabels at=all,
                    yticklabels at=edge left,
                    xlabels at=edge bottom,
                    ylabels at=edge left
                },
                xlabel={\certTimeSymbol\ (pm)},
                ylabel={Simulation},
                xticklabel style={
                    font=\small
                },
                y tick label style={
                    /pgf/number format/fixed, font=\small
                },
                %width=0.45*\textwidth,
                scaled y ticks=false,
                scaled x ticks=false,
                xbar, % to draw a horizontal histogram
                ytick=data,
                nodes near coords, % add data right of the bar
                enlarge x limits={upper, value=0.18}, % to have enough room for the nodes near coords, and just at the right-hand side
                %bar width=4pt,
                %xbar=4.5pt
                xmin=0,
        ]

            \nextgroupplot[title={\normalsize \textbf{Configuration 1}}]
            \addplot[draw=black, fill=gray!50, bar width=4.5pt] table[x index=0, y expr=\coordindex+1]{\data};
            \addplot[draw=black, fill=white, bar width=4.5pt] table[x index=1, y expr=\coordindex+1]{\data};

            \coordinate (top1) at (rel axis cs:0,1);% coordinate at top of the first plot

            \nextgroupplot[title={\normalsize \textbf{Configuration 2}},
                    legend style={
                        at={($(0,0)+(1cm,1cm)$)},
                        legend columns=2,
                        fill=none,
                        draw=none,
                        anchor=south,
                        inner sep=0.2em,
                        column sep=5pt,
                        align=center,
                        %yshift=1ex,
                    },
                    legend to name=leg,
            ]
            \addplot[draw=black, fill=gray!50, bar width=4.5pt] table[x index=2, y expr=\coordindex+1]{\data};
            \addlegendentry{\certTimeSymbol\ (\approachSOTA)}
            \addplot[draw=black, fill=white, bar width=4.5pt] table[x index=3, y expr=\coordindex+1]{\data};
            \addlegendentry{\certTimeSymbol\ (Our appr.)}

            \coordinate (top2) at (rel axis cs:1,1);% coordinate at top of the second plot

        \end{groupplot}

        % legend
        \coordinate (legend pos pre) at ($(top1)!.5!(top2)$);
        \coordinate (legend pos) at (legend pos pre |- current bounding box.north);
        \node[above] at (legend pos)
        {\pgfplotslegendfromname{leg}};

    \end{tikzpicture}
\end{adjustbox}
\end{frame}
